\section{Faces and tight cuts of the perfect matching polytope}
\label{geom:sec-polytope}

The geometric input to the entropy argument is a bound on the number
of independent equations defining a face on its positive coordinates.
We prove the odd-cut description by gluing matching laws across a tight
cut. The same gluing operation then combines directions in the two
contracted polytopes. Both directions preserve the total mass on the cut,
which saves one compatibility equation and yields the sharp bound
$3m-2$. The laminar proof of the same estimate is given in
Appendix~\ref{a:section}.

Throughout this section, \(G=(V,E)\) is a finite undirected loopless
multigraph of even order \(|V|=2m\). Parallel edges have distinct
labels and coordinates. Write \(\mathcal M(G)\) for its perfect
matchings, put \(a(M)=\mathbf 1_M\in\{0,1\}^{E}\), and define
\[
 P(G)=\operatorname{conv}\{a(M):M\in\mathcal M(G)\}.
\]
The convex hull of an empty matching family is empty. If \(V\) is
empty, its unique perfect matching is the empty matching and its
polytope is a singleton in \(\mathbb R^\varnothing\).
For \(S\subseteq V\), let \(\delta(S)\) be the set of edges with
exactly one endpoint in \(S\), and put
\(\chi_S=\mathbf 1_{\delta(S)}\).
We abbreviate \(\delta(\{v\})\) to \(\delta(v)\), and write
\(x(A)=\sum_{e\in A}x_e\) for \(A\subseteq E\).

\subsection{Contraction and the odd-cut description}

For an odd set $S\subsetneq V$, form $G_1$ by retaining $S$ and
contracting $\overline S=V\setminus S$ to a new vertex, discarding edges
internal to $\overline S$. Form $G_2$ by interchanging the two sides.
Every edge of $C=\delta(S)$ keeps its original label in both graphs;
edges that become parallel are not identified. We call the cut
\emph{nontrivial} when both sides have at least three vertices. In that
case the contracted graphs have smaller even orders.

\begin{lemma}[Gluing across an odd cut]
\label{geom:lem-gluing}\label{a:lem-tight-cut-gluing}
Suppose $y_i\in P(G_i)$, for $i=1,2$, and their coordinates agree on
$C$. There is a point $y\in P(G)$ with the internal coordinates of the
two $y_i$ and their common crossing coordinates.
\end{lemma}

\begin{proof}
Choose a matching law of mean $y_i$ in each contracted graph.
A matching uses exactly one edge at the contracted vertex, so the
common coordinates $q_e=(y_1)_e=(y_2)_e$, $e\in C$, sum to one.
Choose a label $e$ with these probabilities, then draw the two
matchings independently from their respective laws conditioned on
using $e$. A label with $q_e=0$ is never chosen and needs no
conditional law.

Remove the contracted vertices and replace the two copies of the
chosen crossing edge by its original edge. The internal matchings
and this edge cover every original vertex once. The crossing
marginals are $q_e$, and the law of total probability preserves every
internal marginal. The resulting matching law therefore has mean $y$.
\end{proof}

The following is the perfect-matching form of Edmonds' theorem
\citep[Section~2, Theorem~(P)]{Edmonds1965}. We include the proof for
loopless multigraphs, with edge labels retained through contractions.

\begin{theorem}[Perfect matching polytope]
\label{geom:thm-polytope}
For every graph \(G\) as above,
\begin{equation}
\label{geom:eq-polytope}
 P(G)=\left\{x\in\mathbb R^E:
 \begin{aligned}
 x_e&\ge0 &&(e\in E),\\
 x(\delta(v))&=1 &&(v\in V),\\
 x(\delta(S))&\ge1 &&(S\subseteq V,\ |S|\text{ odd})
 \end{aligned}\right\}.
\end{equation}
\end{theorem}

\begin{proof}
Denote the set on the right by \(Q\). Every perfect matching has
unit degrees and crosses every odd cut, so \(P(G)\subseteq Q\).
The degree equation at either endpoint of an edge gives
\(0\le x_e\le1\) for every \(x\in Q\). Thus \(Q\) is bounded.
If it is empty there is nothing to prove.

We use the elementary fact that a nonempty bounded polytope is the
convex hull of its extreme points. Indeed, unless a point is extreme,
there is a nonzero direction annihilating all its active constraint
rows. The feasible interval in this direction has two endpoints, and
the point is a convex combination of them. Each endpoint has a new
independent active row: the new row has nonzero directional change,
whereas every old active row annihilates the direction. Repeating
within the resulting lower-dimensional faces terminates at extreme
points.

We show that the extreme points of \(Q\) belong to \(P(G)\), by
induction on the even order. The assertion for order zero is
immediate. At order two, \(Q\) is the simplex
\(\{x\ge0:\sum_{e\in E}x_e=1\}\) on the distinguished parallel
edges, which is exactly their matching polytope and is empty if there
is no edge.

Now let \(|V|\ge4\), and let \(x\) be an extreme point of \(Q\).
Suppose that an odd set \(S\) with
\(3\le |S|\le |V|-3\) has \(x(\delta(S))=1\).
Put $C=\delta(S)$, use the two contracted graphs defined above, and
let $x_i$ be the restriction of $x$ to the surviving edge labels of $G_i$.

Both contracted graphs have smaller even order. The degree equations
at retained vertices are unchanged, and the degree sum at the new
vertex is \(x(C)=1\). To check an odd-cut inequality in a contracted
graph, complement its side if necessary so that it excludes the new
vertex. The complement remains odd because that graph has even
order. Its cut then has exactly the same edge labels and sum as the
cut of the corresponding odd set of original vertices. Thus each
\(x_i\) satisfies the smaller system in
Equation~\eqref{geom:eq-polytope}. By induction, it is the mean of a
law on perfect matchings of \(G_i\).

Their crossing coordinates agree, so Lemma~\ref{geom:lem-gluing}
gives $x\in P(G)$.

It remains to consider an extreme point at which every such
nontrivial odd cut is slack. A nonzero vector supported on the
positive coordinates of \(x\) and preserving all degree equations
would allow sufficiently small feasible perturbations in both signs:
the positive coordinates stay positive, all nontrivial odd cuts
retain positive slack, and singleton cuts and their complements are
controlled by the degree equations. Therefore the columns of the
unsigned vertex-edge incidence matrix on the positive edges are
linearly independent.

Consider a connected component of the positive-edge graph. There are
no isolated vertices, by the degree equations. A vertex incident to
only one positive edge forces that edge to have value one; its other
endpoint then has no other positive incident edge. These vertices
therefore occur in isolated edge components. Every remaining
component has minimum degree at least two. Column independence
implies that its number of edges is at most its number of vertices.
Counting incidences forces equality and degree two at every vertex,
so the component is a cycle, possibly a two-edge parallel cycle. An
even cycle gives a nonzero alternating-sign dependence among the
incidence columns. An odd cycle component has cut sum zero on its
odd vertex set, contrary to feasibility. Thus only isolated
weight-one edge components remain, and \(x\) is a perfect matching
vector. This completes the induction.
\end{proof}

\subsection{Directions in the minimal face}

For \(x\in P(G)\), define
\[
 T(x)=\{h\in\mathbb R^E:
     x+th\in P(G)\text{ for all sufficiently small }|t|\}.
\]
Let \(F_x\) be the face obtained by keeping equality in every
inequality of Equation~\eqref{geom:eq-polytope} that is tight at
\(x\). Define
\[
 \mathcal A_x=\{a(M):M\in\mathcal M(G),\ a(M)\in F_x\}.
\]
The next lemma identifies \(T(x)\) as a linear space and explains
why this face is the one relevant to laws with mean \(x\).

\begin{lemma}[Minimal face and its directions]
\label{geom:lem-face}
For every \(x\in P(G)\), the space \(T(x)\) consists exactly of
the vectors satisfying
\begin{equation}
\label{geom:eq-directions}
\begin{aligned}
 h_e&=0 &&\text{if }x_e=0,\\
 h(\delta(v))&=0 &&(v\in V),\\
 h(\delta(S))&=0
   &&\text{if }|S|\text{ is odd and }x(\delta(S))=1.
\end{aligned}
\end{equation}
Moreover, \(F_x\) is the minimal face containing \(x\), and
\begin{equation}
\label{geom:eq-face-span}
 \operatorname{aff}F_x=x+T(x),\qquad
 F_x=\operatorname{conv}\mathcal A_x,\qquad
 T(x)=\operatorname{span}\{a-x:a\in\mathcal A_x\}.
\end{equation}
The point \(x\) lies in the relative interior of \(F_x\), and every
law on perfect matching vectors with mean \(x\) is supported on
\(\mathcal A_x\).
\end{lemma}

\begin{proof}
Feasibility in both signs forces every active constraint row to
annihilate the direction. Conversely, the equations in
Equation~\eqref{geom:eq-directions} preserve all active constraints,
and sufficiently small perturbations in both signs preserve every
inactive inequality, since there are finitely many and each has
positive slack. This proves the direction characterization.

The same argument gives a relative neighborhood of \(x\) in
\(x+T(x)\) contained in \(F_x\). Explicitly, for each inactive
inequality \(b_i\cdot y\ge c_i\) with nonzero normal, choose a
radius smaller than its slack at \(x\) divided by \(2\|b_i\|_2\),
and take the minimum of these finitely many positive radii. If no
such row exists, any positive radius suffices. All active rows stay
equalities on \(x+T(x)\), and zero-normal rows are constant.
Since \(F_x\subseteq x+T(x)\), this proves its affine hull and
relative interiority.

For any \(y\in F_x\), a short extension from \(y\) through \(x\)
remains in this relative neighborhood. Thus \(x\) lies in the
relative interior of a segment in \(F_x\) having \(y\) as an
endpoint. Every face containing \(x\) must contain \(y\), proving
minimality.

Finally, express a point of \(F_x\) as a convex combination of
matching vectors. Each inequality defining \(F_x\) has nonnegative
slack at every vector and zero average slack. Every positive-weight
vector must therefore belong to \(F_x\). This proves the convex-hull
and support assertions. Taking the direction space of this affine
hull gives the final equality in
Equation~\eqref{geom:eq-face-span}.
\end{proof}

Write \(d(x)=\dim F_x=\dim T(x)\). For a point whose coordinates
are all positive, let
\[
 \mathcal S_x=\{S\subseteq V:|S|\text{ is odd},\
                                x(\delta(S))=1\}.
\]
There are no active nonnegativity constraints, and the degree rows
are already the cut rows of singletons. Lemma~\ref{geom:lem-face}
therefore gives the exact identity
\begin{equation}
\label{geom:eq-rank}
 |E|-d(x)=\dim\operatorname{span}\{\chi_S:S\in\mathcal S_x\}.
\end{equation}
We now bound the rank on the right.

\subsection{Sharp codimension from compatible directions}

The gluing lemma controls dimensions as well as feasible means.
Directions in the two contracted faces can be glued whenever their cut
coordinates agree. Their cut-coordinate sums are already zero, so this
agreement imposes at most $|C|-1$ independent conditions.

\begin{proposition}[Codimension on positive coordinates]
\label{geom:prop-rank}
Let $|V|=2m\ge2$, and suppose $x\in P(G)$ has $x_e>0$ for every
edge. Then
\begin{equation}
\label{geom:eq-codimension}
 |E|-d(x)\le3m-2.
\end{equation}
\end{proposition}

\begin{proof}
We induct on the even order $n=2m$. At order two, the polytope is the
simplex on the distinguished parallel edges. A positive point has
dimension $|E|-1$, giving codimension one as required.

Let $n\ge4$. If no nontrivial odd cut is tight, positivity and
Lemma~\ref{geom:lem-face} show that only the $n$ degree equations
restrict $T(x)$. Thus $|E|-d(x)\le n\le3n/2-2$.
Otherwise choose a tight nontrivial cut $C=\delta(S)$, and let $x_i$
be the restricted means in the two contracted graphs $G_i$.
The feasibility argument in Theorem~\ref{geom:thm-polytope} gives
$x_i\in P(G_i)$, and every retained coordinate is positive. Put
\[
 n_i=|V(G_i)|,\qquad s_i=|E(G_i)|,\qquad d_i=\dim T(x_i).
\]
Each crossing label appears in both graphs, whereas each internal
label appears in one. Consequently
\begin{equation}\label{geom:eq-contraction-sizes}
 n_1+n_2=n+2,\qquad s_1+s_2=|E|+|C|,\qquad 4\le n_i<n.
\end{equation}

Consider the space of compatible directions
\[
 W=\{(h_1,h_2)\in T(x_1)\oplus T(x_2):
                         (h_1)_e=(h_2)_e\text{ for }e\in C\}.
\]
Each $h_i$ preserves the degree at its contracted vertex, so
$\sum_{e\in C}(h_i)_e=0$. The linear map
$(h_1,h_2)\mapsto h_1|_C-h_2|_C$ therefore takes values in
\[
 \{z\in\mathbb R^C:\textstyle\sum_{e\in C}z_e=0\}.
\]
Since $x(C)=1$, the set $C$ is nonempty and this target has dimension
$|C|-1$. Its kernel is $W$, whence
\[
                         \dim W\ge d_1+d_2-|C|+1.
\]

Map $(h_1,h_2)\in W$ to the vector $h$ on the original edges with
their internal coordinates and common crossing coordinates. This map
is injective: restricting $h$ recovers both $h_i$. For any such pair,
choose $\varepsilon>0$ so that $x_i+th_i\in P(G_i)$ for both $i$
and every $|t|<\varepsilon$. Their crossing coordinates agree for
each $t$, so Lemma~\ref{geom:lem-gluing} yields $x+th\in P(G)$.
Thus $h\in T(x)$ and $d(x)\ge\dim W$. In particular every parent
cut tight at $x$ is preserved, including cuts crossing $S$, because
both signs of the perturbation are feasible.

Using Equation~\eqref{geom:eq-contraction-sizes} and induction gives
\begin{align*}
 |E|-d(x)
 &\le (s_1-d_1)+(s_2-d_2)-1\\
 &\le (3n_1/2-2)+(3n_2/2-2)-1
   =3n/2-2.
\end{align*}
The saving of one compatibility condition is exactly what closes
this recurrence at the stated constant.
\end{proof}

No connectedness assumption or strict upper bound $x_e<1$ was used.
Section~\ref{ex:sec-triangle-expansions} constructs positive means
attaining this codimension at every positive even order.

For an arbitrary mean $x\in P(G)$ on $2m\ge2$ vertices, deleting just
its zero-coordinate edges leaves the same vertex set and identifies the supported minimal face
with the minimal face in the resulting positive-edge graph. Indeed,
every law at $x$ omits those edges, and coordinate restriction and zero
extension are inverse on the corresponding face. Hence the supported
form of the estimate is
\begin{equation}\label{geom:eq-supported-rank}
                |\supp x|-d(x)\le3m-2.
\end{equation}
The next reduction additionally removes coordinates equal to one.

\subsection{Removing fixed coordinates}

We finally transfer the rank estimate to an arbitrary polytope
point. This reduction will also allow entropy comparisons to be
made entirely on the fractional coordinates.
For \(x\in P(G)\), set
\[
 E_f(x)=\{e\in E:0<x_e<1\},\qquad
 V_f(x)=\{v\in V:v\text{ meets an edge of }E_f(x)\}.
\]

\begin{lemma}[Fractional coordinates and the minimal face]
\label{geom:lem-reduction}
Let \(x\in P(G)\). The number of vertices in \(V_f(x)\) is even;
write \(|V_f(x)|=2m_f\). Let
\(G_f=(V_f(x),E_f(x))\) and \(\bar x=x|_{E_f(x)}\).
The coordinate extension
\begin{equation}
\label{geom:eq-extension}
 \iota(z)_e=
 \begin{cases}
 z_e,&e\in E_f(x),\\
 x_e,&e\notin E_f(x)
 \end{cases}
\end{equation}
is an affine isomorphism from \(P(G_f)\) onto the face
\[
 Q_x=\{y\in P(G):y_e=x_e\text{ for }e\notin E_f(x)\}.
\]
It sends the minimal face of \(\bar x\) onto \(F_x\). For every
\(z\in P(G_f)\), removing the fixed edges gives a bijection between
laws with mean \(\iota(z)\) and laws on perfect matchings of \(G_f\)
with mean \(z\), preserving their probabilities. Moreover,
\begin{equation}
\label{geom:eq-fractional-rank}
 \sum_{e\in E_f(x)}x_e=m_f,\qquad
 d(x)\ge |E_f(x)|-3m_f+2\quad\text{if }m_f\ge1.
\end{equation}
If \(m_f=0\), then \(x\) is a matching vector and \(d(x)=0\).
\end{lemma}

\begin{proof}
An edge with coordinate one exhausts the degree sum at both
endpoints, so neither endpoint meets a fractional edge. At a vertex
outside \(V_f(x)\), all incident coordinates are zero or one and
sum to one; exactly one incident edge therefore has coordinate one.
These fixed edges form a perfect matching of \(V\setminus V_f(x)\),
so \(|V_f(x)|\) is even.

Every law of mean \(x\) uses each value-one edge and omits each
value-zero edge with probability one. Deleting the fixed edges gives
a perfect matching law on \(G_f\) with mean \(\bar x\), proving
\(\bar x\in P(G_f)\). Conversely, adjoining the fixed edges to any
perfect matching of \(G_f\) gives a perfect matching of \(G\).

The inequalities \(0\le y_e\le1\) hold throughout \(P(G)\), so
requiring any of these coordinates to equal zero or one defines a
face; their intersection \(Q_x\) is a face containing \(x\).
The preceding matching extension shows
\(\iota(P(G_f))\subseteq Q_x\). For the reverse inclusion, express
any point of \(Q_x\) as a law on perfect matchings. Its fixed zero
and one marginals force every positive-probability matching to omit
the zero edges and contain the fixed edges. Removing those edges
leaves a perfect matching of \(G_f\). This proves both equality of
the image with \(Q_x\) and the stated bijection of laws for every
mean. Coordinate restriction is the affine inverse of \(\iota\).

An affine isomorphism preserves faces and their relative interiors.
The image of the minimal face at \(\bar x\) is thus a face of
\(Q_x\), and hence of \(P(G)\), containing \(x\) in its relative
interior. By Lemma~\ref{geom:lem-face}, it is exactly \(F_x\).
In particular, their dimensions are equal. One can also see this
directly on directions: a two-sided direction at \(x\) vanishes on
every fixed zero or one coordinate, and its restriction to
\(E_f(x)\) is a two-sided direction at \(\bar x\); extension gives
the inverse correspondence.

If \(m_f\ge1\), every coordinate of \(\bar x\) is positive, and
Proposition~\ref{geom:prop-rank} applied to \(G_f\) gives
\[
 d(x)=\dim F_{\bar x}\ge |E_f(x)|-3m_f+2.
\]
If \(m_f=0\), all coordinates of \(x\) are integral, and the
degree equations identify \(x\) as a single matching vector.
Finally, at every vertex of \(V_f(x)\) the fractional incident
coordinates sum to one. Summing those degree equations counts each
fractional edge twice and gives
\(\sum_{e\in E_f(x)}x_e=m_f\).
\end{proof}
